\section{Заключение}

В ходе выполнения лабораторной работы были изучены принципы мониторинга, профилирования и оптимизации производительности программных продуктов на платформе Java. На примере веб-приложения и утилиты автоматизации тестирования были практически применены технологии JMX \textit{(Java Management Extensions)}, а также инструменты для анализа состояния виртуальной машины Java.

Были решены следующие задачи:
\begin{itemize}
  \item проектирование и программная реализация управляемых компонентов MBean (\texttt{PointCounter} и \texttt{MissPercentage}) для динамического сбора метрик пользовательской активности, а также интеграция асинхронного механизма JMX-уведомлений;
  \item проведение мониторинга приложения через утилиту \texttt{JConsole} с получением данных используемой JVM, валидацией измеряемых атрибутов и ручным перехватом генерируемых событий;
  \item профилирование распределения памяти в среде \texttt{VisualVM} с целью нахождения наиболее ресурсоемких системных классов и пользовательских сущностей, которые формируют общую нагрузку на кучу;
  \item локализация и анализ архитектурных задержек в главном потоке выполнения с помощью встроенного профайлера методов;
  \item обнаружение, диагностика и устранение утечки памяти, вызванной неконтролируемым накоплением логов в статической коллекции библиотеки \texttt{HttpUnit}.
\end{itemize}

Полученные навыки анализа дампов кучи, трассировки потоков и управления жизненным циклом ресурсов --- основа для обеспечения надежности, отказоустойчивости и высокой пропускной способности современных систем на базе Java.